\chapter{The Schauder Approach}

Throughout, $L=a^{ij}D_{ij}+b^iD_i+c$ is a linear second-order operator and
$d_x=\operatorname{dist}(x,\partial\Omega)$, $d_{x,y}=\min(d_x,d_y)$.  The
interior and boundary weighted H\"older seminorms are those of (4.17)--(4.18).

\begin{lemma}[Constant-coefficient reduction]
\label{gt-ch06-lem-6-1}
\dcref{ch6:6.1}
\group{gt-ch06}
Let $A=[A^{ij}]$ be constant and satisfy
\[
\lambda|\xi|^2\le A^{ij}\xi_i\xi_j\le\Lambda|\xi|^2\qquad(\xi\in\mathbb R^n),
\]
with $0<\lambda\le\Lambda$.  If $L_0u=A^{ij}D_{ij}u=f$ in an open set
$\Omega$, then
\begin{equation}\tag{6.4}
|u|_{2,\alpha;\Omega}^*\le C\bigl(|u|_{0;\Omega}+|f|_{0,\alpha;\Omega}^{(2)}\bigr).
\end{equation}
If $\Omega\subset\mathbb R^n_+$ has a boundary portion $T\subset\{x_n=0\}$,
$u=0$ on $T$, and $u,f$ extend with the stated H\"older regularity to $T$, then
\begin{equation}\tag{6.5}
|u|_{2,\alpha;\Omega\cup T}^*\le C\bigl(|u|_{0;\Omega}+|f|_{0,\alpha;\Omega\cup T}^{(2)}\bigr).
\end{equation}
Here and below $C=C(n,\alpha,\lambda,\Lambda)$ unless additional dependence is
displayed.
\end{lemma}
\begin{proof}
\sketch An orthogonal diagonalization of $A$ followed by the scaling
$y=xQ$, $Q=P\operatorname{diag}(\lambda_i^{-1/2})$, changes $L_0$ into the
Laplacian.  The map satisfies
$\Lambda^{-1/2}|x|\le|xQ|\le\lambda^{-1/2}|x|$; hence all weighted interior
and flat-boundary norms are equivalent:
\begin{equation}
\tag{6.7}
C^{-1}|v|^*_{k,\alpha;\Omega\cup T}
\leq|v\circ Q^{-1}|^*_{k,\alpha;Q(\Omega\cup T)}
\leq C|v|^*_{k,\alpha;\Omega\cup T}.
\end{equation}
The constants depend only on
$n,\lambda,\Lambda$ and the order.  Apply the Poisson estimates in the
transformed domain and transfer them back.
\end{proof}

\section*{6.1. The Schauder Interior Estimates}

For $u\in C^{2,\alpha}(\Omega)$, $\varepsilon>0$, and
$j+\beta<2+\alpha$,
\begin{equation}\tag{6.8}
[u]^*_{j;\beta;\Omega}\le C(\varepsilon)|u|_{0;\Omega}+\varepsilon [u]^*_{2,\alpha;\Omega},
\qquad 0\le j\le2,
\end{equation}
and
\begin{equation}\tag{6.9}
|u|^*_{j;\beta;\Omega}\le C(\varepsilon)|u|_{0;\Omega}+\varepsilon [u]^*_{2,\alpha;\Omega}.
\end{equation}

For a real $\sigma$ and integer $k\ge0$, set
\[
[f]_{k;\Omega}^{(\sigma)}=\sup_{x\in\Omega,|\beta|=k}d_x^{k+\sigma}|D^\beta f(x)|,
\]
\[
[f]_{k,\alpha;\Omega}^{(\sigma)}=
\sup_{x\ne y,|\beta|=k}d_{x,y}^{k+\alpha+\sigma}
\frac{|D^\beta f(x)-D^\beta f(y)|}{|x-y|^\alpha},
\]
\begin{equation}
\tag{6.10}
|f|_{k;\Omega}^{(\sigma)}=\sum_{j=0}^k[f]_{j;\Omega}^{(\sigma)},
\end{equation}
and $|f|_{k,\alpha;\Omega}^{(\sigma)}=|f|_{k;\Omega}^{(\sigma)}
+[f]_{k,\alpha;\Omega}^{(\sigma)}$.
For $\sigma=0$ these are the starred norms.  Products satisfy
\begin{equation}\tag{6.11}
|fg|_{0;\Omega}^{(\sigma+\tau)}
\le |f|_{0;\Omega}^{(\sigma)}|g|_{0;\Omega}^{(\tau)}\qquad(\sigma+\tau\ge0).
\end{equation}

\begin{theorem}[Interior Schauder estimate]
\label{gt-ch06-thm-6-2}
\source{gilbarg-trudinger:page-0103}
\source{gilbarg-trudinger:page-0104}
\source{gilbarg-trudinger:page-0105}
\dcref{ch6:6.2}
\group{gt-ch06}

Let $\Omega\subset\mathbb R^n$ be open, $u\in C^{2,\alpha}(\Omega)$ bounded, and
\[
Lu=a^{ij}D_{ij}u+b^iD_i u+cu=f,
\]
where $f\in C^\alpha(\Omega)$ and
\begin{equation}\tag{6.12}
a^{ij}(x)\xi_i\xi_j\ge\lambda|\xi|^2,
\end{equation}
\begin{equation}\tag{6.13}
|a^{ij}|_{0;\Omega}^{(0)}+|b^i|_{0;\Omega}^{(1)}+|c|_{0;\Omega}^{(2)}\le\Lambda.
\end{equation}
Then
\begin{equation}\tag{6.14}
|u|_{2,\alpha;\Omega}^*\le C\bigl(|u|_{0;\Omega}+|f|_{0,\alpha;\Omega}^{(2)}\bigr).
\end{equation}
\end{theorem}
\begin{proof}
\sketch Interpolation reduces the claim to the starred second-derivative
H\"older seminorm.  Exhaust $\Omega$ by relatively compact open sets and pass to
the limit, so the seminorm may be assumed finite.  For $x_0,y_0$ with
$d_{x_0}=d_{x_0,y_0}$, put $d=\mu d_{x_0}$ and freeze the principal
coefficients at $x_0$:
\begin{equation}\tag{6.15}
a^{ij}(x_0)D_{ij}u=F=(a^{ij}(x_0)-a^{ij}(x))D_{ij}u-b^iD_i u-cu+f.
\end{equation}
Apply the constant-coefficient estimate in $B_d(x_0)$ and distinguish
$|x_0-y_0|\le d/2$ from $|x_0-y_0|>d/2$. This gives
\begin{equation}
\tag{6.16}
d_{x_0,y_0}^{2+\alpha}\frac{|D^2u(x_0)-D^2u(y_0)|}{|x_0-y_0|^\alpha}
\le C\mu^\alpha [u]^*_{2,\alpha;\Omega}
+C(\mu)\bigl(|u|_{0;\Omega}+|F|_{0,\alpha;B}^{(2)}\bigr).
\end{equation}
The frozen right-hand side satisfies
\begin{equation}
\tag{6.17}
|F|_{0,\alpha;B}^{(2)}
\leq\sum_{i,j}|(a^{ij}(x_0)-a^{ij})D_{ij}u|_{0,\alpha;B}^{(2)}
+\sum_i|b^iD_i u|_{0,\alpha;B}^{(2)}+|cu|_{0,\alpha;B}^{(2)}
+|f|_{0,\alpha;B}^{(2)}.
\end{equation}
For $g\in C^\alpha(\Omega)$ and $\mu\le1/2$,
\begin{equation}
\tag{6.18}
|g|_{0,\alpha;B}^{(2)}
\leq4\mu^2|g|_{0;\Omega}^{(2)}
+8\mu^{2+\alpha}[g]_{0,\alpha;\Omega}^{(2)}
\leq8\mu^2|g|_{0,\alpha;\Omega}^{(2)}.
\end{equation}
Using (6.8), (6.9), (6.11), and (6.18), the three coefficient terms obey
\begin{equation}
\tag{6.19}
\sum_{i,j}|(a^{ij}(x_0)-a^{ij})D_{ij}u|_{0,\alpha;B}^{(2)}
\leq32n^2\Lambda\mu^{2+\alpha}
\bigl(C(\mu)|u|_{0;\Omega}+2\mu^\alpha[u]^*_{2,\alpha;\Omega}\bigr),
\end{equation}
\begin{equation}
\tag{6.20}
\sum_i|b^iD_i u|_{0,\alpha;B}^{(2)}
\leq8n\Lambda\mu^2C(\mu)|u|_{0;\Omega}
+\mu^{2\alpha}[u]^*_{2,\alpha;\Omega},
\end{equation}
and
\begin{equation}
\tag{6.21}
|cu|_{0,\alpha;B}^{(2)}
\leq8\Lambda\mu^2C(\mu)|u|_{0;\Omega}
+\mu^{2\alpha}[u]^*_{2,\alpha;\Omega}.
\end{equation}
Finally,
\begin{equation}
\tag{6.22}
|f|_{0,\alpha;B}^{(2)}
\leq8\mu^2|f|_{0,\alpha;\Omega}^{(2)}.
\end{equation}
Choose $\mu$ so that $C\mu^\alpha\le1/2$ and take the supremum.
\end{proof}

\begin{corollary}[Local interior Schauder estimate and compactness]
\label{gt-ch06-cor-6-3}
\dcref{ch6:6.3}
\group{gt-ch06}
If $\Omega$ is bounded, $u\in C^{2,\alpha}(\Omega)$ solves $Lu=f$, the
coefficients and $f$ are $C^\alpha$, and $\Omega'\Subset\Omega$ has distance
at least $d$ from $\partial\Omega$, then
\begin{equation}\tag{6.23}
d|Du|_{0;\Omega'}+d^2|D^2u|_{0;\Omega'}+d^{2+\alpha}[D^2u]_{\alpha;\Omega'}
\le C\bigl(|u|_{0;\Omega}+|f|_{0,\alpha;\Omega}\bigr).
\end{equation}
Consequently uniformly bounded families with locally uniform ellipticity and
H\"older coefficient bounds are equicontinuous together with their first and
second derivatives on compact subsets.
\end{corollary}

\section*{6.2. Boundary and Global Estimates}

A bounded domain is of class $C^{k,\alpha}$ if each boundary point has a
$C^{k,\alpha}$ diffeomorphism onto a set whose domain and boundary portions
are mapped to $\mathbb R^n_+$ and $\{x_n=0\}$, respectively.  A boundary
portion is $C^{k,\alpha}$ when this holds locally along that portion.  These
classes are nested whenever $j+\beta<k+\alpha$.

\begin{lemma}[Flat-boundary Schauder estimate]
\label{gt-ch06-lem-6-4}
\dcref{ch6:6.4}
\group{gt-ch06}
Let $\Omega\subset\mathbb R^n_+$ have a flat boundary portion $T\subset\{x_n=0\}$.
If $u\in C^{2,\alpha}(\Omega\cup T)$ solves $Lu=f$ in $\Omega$, $u=0$ on
$T$, and
\begin{equation}\tag{6.26}
|a^{ij}|_{0,\alpha;\Omega\cup T}^{(0)},\ |b^i|_{0,\alpha;\Omega\cup T}^{(1)},\ |c|_{0,\alpha;\Omega\cup T}^{(2)}\le\Lambda,
\quad |f|_{0,\alpha;\Omega\cup T}^{(2)}<\infty,
\end{equation}
then
\begin{equation}\tag{6.27}
|u|^*_{2,\alpha;\Omega\cup T}
\le C\bigl(|u|_{0;\Omega}+|f|_{0,\alpha;\Omega\cup T}^{(2)}\bigr).
\end{equation}
\end{lemma}
\begin{proof}
\sketch Repeat the freezing argument for Theorem 6.2 using the flat-boundary
part of the constant-coefficient lemma and the boundary versions of (6.8)--(6.9).
\end{proof}

For a flat boundary portion put
\[
\bar d_x=\operatorname{dist}(x,\partial\Omega-T),\qquad
\bar d_{x,y}=\min(\bar d_x,\bar d_y),
\]
and define starred boundary seminorms by replacing $d$ with $\bar d$ in the
interior formulas; in particular
\[
[u]^*_{k,\alpha;\Omega\cup T}=
\sup_{x\ne y,|\beta|=k}\bar d_{x,y}^{k+\alpha}
\frac{|D^\beta u(x)-D^\beta u(y)|}{|x-y|^\alpha}.
\]
The boundary interpolation inequalities are
\begin{equation}
\tag{6.24}
[u]^*_{j,\beta;\Omega\cup T}
\leq C(\varepsilon)|u|_{0;\Omega}
+\varepsilon[u]^*_{2,\alpha;\Omega\cup T},
\end{equation}
\begin{equation}
\tag{6.25}
|u|^*_{j,\beta;\Omega\cup T}
\leq C(\varepsilon)|u|_{0;\Omega}
+\varepsilon[u]^*_{2,\alpha;\Omega\cup T},
\end{equation}
for $j+\beta<2+\alpha$.
If $\psi$ is a $C^{k,\alpha}$ diffeomorphism with bounded inverse on a compact
set, then for $x'=\psi(x)$,
\begin{equation}\tag{6.29}
K^{-1}|x-y|\le|x'-y'|\le K|x-y|,
\end{equation}
and the corresponding interior, boundary, and weighted norms are equivalent:
\begin{equation}\tag{6.30}
K^{-1}|u|^*_{j,\beta;\Omega\cup T}
\le|u\circ\psi^{-1}|^*_{j,\beta;\Omega'\cup T'}
\le K|u|^*_{j,\beta;\Omega\cup T}.
\end{equation}

\begin{lemma}[Schauder estimate near a boundary point]
\label{gt-ch06-lem-6-5}
\dcref{ch6:6.5}
\group{gt-ch06}
Let $\Omega$ be a $C^{2,\alpha}$ domain and $u\in C^{2,\alpha}(\bar\Omega)$ solve
$Lu=f$ with $u=0$ on $\partial\Omega$, $f\in C^\alpha(\bar\Omega)$, and
\begin{equation}\tag{6.31}
|a^{ij}|_{0;\Omega},\ |b^i|_{0;\Omega},\ |c|_{0;\Omega}\le\Lambda.
\end{equation}
There are $\delta>0$ and $C=C(n,\alpha,\lambda,\Lambda,\Omega)$ such that
for every $x_0\in\partial\Omega$,
\begin{equation}\tag{6.32}
|u|_{2,\alpha;B_\delta(x_0)\cap\Omega}
\le C\bigl(|u|_{0;\Omega}+|f|_{0,\alpha;\Omega}\bigr).
\end{equation}
\end{lemma}
\begin{proof}
\sketch Straighten the boundary by $y=\psi(x)$. The transformed operator has
ellipticity and coefficient bounds
\begin{equation}
\tag{6.33}
\widetilde\lambda=\lambda/K,
\end{equation}
\begin{equation}
\tag{6.34}
|\widetilde a^{ij}|_{0,\alpha;D'},\
|\widetilde b^i|_{0,\alpha;D'},\
|\widetilde c|_{0,\alpha;D'}
\leq\widetilde\Lambda=K\Lambda,
\qquad |\widetilde f|_{0,\alpha;D'}<\infty.
\end{equation}
Lemma 6.4 and (6.30) give, on a smaller boundary ball,
\begin{equation}
\tag{6.35}
|u|_{2,\alpha;B''}
\leq C\bigl(|u|_{0;\Omega}+|f|_{0,\alpha;\Omega}\bigr).
\end{equation}
A finite boundary covering gives a uniform $\delta$.
\end{proof}

\begin{theorem}[Global Schauder estimate]
\label{gt-ch06-thm-6-6}
\dcref{ch6:6.6}
\group{gt-ch06}
Let $\Omega$ be $C^{2,\alpha}$, $u\in C^{2,\alpha}(\bar\Omega)$ solve $Lu=f$,
$f\in C^\alpha(\bar\Omega)$, and let the coefficients satisfy (6.12) and
\[
|a^{ij}|_{0,\alpha;\Omega},\ |b^i|_{0,\alpha;\Omega},\ |c|_{0,\alpha;\Omega}\le\Lambda.
\]
If $u=\varphi$ on $\partial\Omega$ with $\varphi\in C^{2,\alpha}(\bar\Omega)$,
then
\begin{equation}\tag{6.36}
|u|_{2,\alpha;\Omega}\le C\bigl(|u|_{0;\Omega}+|\varphi|_{2,\alpha;\Omega}+|f|_{0,\alpha;\Omega}\bigr).
\end{equation}
\end{theorem}
\begin{proof}
\sketch Subtract an extension of $\varphi$ to reduce to zero boundary data.
Use Lemma 6.5 near the boundary and Corollary 6.3 on the interior strip.  For
every $x\in\Omega$ these estimates first give
\begin{equation}
\tag{6.37}
|Du(x)|+|D^2u(x)|
\leq C\bigl(|u|_0+|f|_{0,\alpha}\bigr).
\end{equation}
For
two points in different local balls, bound the H\"older quotient by the global
second-derivative estimate and their separation; the three cases yield (6.36).
\end{proof}

\begin{corollary}[Local Schauder estimate on a boundary portion]
\label{gt-ch06-cor-6-7}
\dcref{ch6:6.7}
\group{gt-ch06}
If $T$ is a $C^{2,\alpha}$ boundary portion, $u\in C^{2,\alpha}(\Omega\cup T)$,
$u=\varphi$ on $T$, and the hypotheses of Theorem 6.6 hold, then for any
$x_0\in T$ and $B=B_\rho(x_0)$ with $\rho<\operatorname{dist}(x_0,\partial\Omega-T)$,
\begin{equation}\tag{6.38}
|u|_{2,\alpha;B\cap\Omega}\le C\bigl(|u|_{0;\Omega}+|\varphi|_{2,\alpha;\Omega}+|f|_{0,\alpha;\Omega}\bigr).
\end{equation}
\end{corollary}

\section*{6.3. The Dirichlet Problem}

\begin{theorem}[Dirichlet solvability by continuity]
\label{gt-ch06-thm-6-8}
\dcref{ch6:6.8}
\group{gt-ch06}
Let $\Omega$ be a $C^{2,\alpha}$ domain and $L$ strictly elliptic with
$C^\alpha(\bar\Omega)$ coefficients and $c\le0$.  If the Poisson Dirichlet
problem has a $C^{2,\alpha}(\bar\Omega)$ solution for all
$f\in C^\alpha(\bar\Omega)$ and $\varphi\in C^{2,\alpha}(\bar\Omega)$, then
\begin{equation}\tag{6.39}
Lu=f\ \text{in }\Omega,\qquad u=\varphi\ \text{on }\partial\Omega
\end{equation}
has a unique solution in $C^{2,\alpha}(\bar\Omega)$ for all such data.
\end{theorem}
\begin{proof}
\sketch Reduce to $\varphi=0$ and use $L_t=tL+(1-t)\Delta$.  Uniform ellipticity,
with
\begin{equation}
\tag{6.40}
\lambda|\xi|^2\leq a^{ij}(x)\xi_i\xi_j
\qquad(x\in\Omega,\ \xi\in\mathbb R^n),
\end{equation}
and use the family
\begin{equation}
\tag{6.41}
L_tu=tLu+(1-t)\Delta u=f,\qquad0\leq t\leq1.
\end{equation}
The maximum-principle $C^0$ bound and (6.36) give
\begin{equation}
\tag{6.42}
|u_t|_{2,\alpha}\leq C|f|_{0,\alpha},
\end{equation}
uniformly in $t$. The method of continuity from
$L_0=\Delta$ proves invertibility for $t=1$.
\end{proof}

\begin{corollary}[Dirichlet solvability on a ball]
\label{gt-ch06-cor-6-9}
\dcref{ch6:6.9}
\group{gt-ch06}
Under Theorem 6.8 with $\Omega=B$ a ball, the Dirichlet problem has a unique
$C^{2,\alpha}(\bar B)$ solution for $f\in C^\alpha(\bar B)$ and
$\varphi\in C^{2,\alpha}(\bar B)$.
\end{corollary}

\begin{lemma}[Partial boundary regularity on a ball]
\label{gt-ch06-lem-6-10}
\dcref{ch6:6.10}
\group{gt-ch06}
Let $T\subset\partial B$ be any portion of a ball, $\varphi\in C^0(\partial B)\cap
C^{2,\alpha}(T)$, and let $L,f$ satisfy Theorem 6.8.  Then the Dirichlet
problem in $B$ has a unique solution
$u\in C^0(\bar B)\cap C^{2,\alpha}(B\cup T)$.
\end{lemma}
\begin{proof}
\sketch Extend $\varphi$ near $T$ with uniform $C^{2,\alpha}$ bounds, approximate
by smooth boundary data $\varphi_k$ satisfying
\begin{equation}
\tag{6.43}
|\varphi_k-\varphi|_{0;B}\to0,
\qquad |\varphi_k|_{2,\alpha;G}\leq C|\varphi|_{2,\alpha;G}.
\end{equation}
Solve by Corollary 6.9 and pass to the uniform and
interior Schauder limits.  Uniqueness follows from the maximum principle.
\end{proof}

For continuous $u$, call $u$ a subsolution (supersolution) of $Lu=f$ if it is
below (above) every solution with the same boundary comparison on each ball
$B\Subset\Omega$.  A $C^2$ function is a subsolution exactly when $Lu\ge f$;
the dual statement holds for supersolutions.  Subsolutions are stable under
replacement by the solution on a compact ball and under finite maxima; the
corresponding statements hold for supersolutions.  If a subsolution is below a
supersolution on $\partial\Omega$ and $c\le0$, either the inequality is strict
inside or the two functions coincide.

For bounded boundary data $\varphi$, a subfunction (superfunction) is a
subsolution (supersolution) below (above) $\varphi$ on $\partial\Omega$.  If
$\Omega$ lies in $0<x_1<d$, then for sufficiently large $\gamma$,
\begin{equation}
\tag{6.44}
v^\pm=\pm\sup|\varphi|
\pm(e^{\gamma d}-e^{\gamma x_1})\frac{\sup|f|}{\lambda}
\end{equation}
are a superfunction and subfunction. Thus $S_\varphi$ is nonempty and bounded
above; define
\[
u(x)=\sup_{v\in S_\varphi}v(x).
\]

\begin{theorem}[Perron solution]
\label{gt-ch06-thm-6-11}
\dcref{ch6:6.11}
\group{gt-ch06}
If $c\le0$, the coefficients and $f$ are locally $C^\alpha$, and the above
supremum is bounded, then $u\in C^{2,\alpha}(\Omega)$ and $Lu=f$ in $\Omega$.
\end{theorem}
\begin{proof}
\sketch Solve on balls, use stability of subsolutions and the interior compactness
from Corollary 6.3, and apply the strong maximum principle.
\end{proof}

At a boundary point $x_0$ where $\varphi$ is continuous, an upper (lower)
barrier is a sequence of superfunctions (subfunctions) $w_i^\pm$ with
$w_i^\pm(x_0)\to\varphi(x_0)$.  A barrier means both sequences.

\begin{lemma}[Boundary convergence from a barrier]
\label{gt-ch06-lem-6-12}
\dcref{ch6:6.12}
\group{gt-ch06}
If the Perron solution of Theorem 6.11 has a barrier at $x_0$, then
$u(x)\to\varphi(x_0)$ as $x\to x_0$.
\end{lemma}
\begin{proof}
\sketch Every subfunction is bounded above by every superfunction, hence
$w_i^-\le u\le w_i^+$.  Letting $x\to x_0$ and then $i\to\infty$ gives the claim.
\end{proof}

If $c=f=0$, a nonnegative supersolution $w$ with $w(x_0)=0$ and
$w>0$ on $\partial\Omega\setminus\{x_0\}$ determines a barrier.  More generally,
if $Lw\le-1$ with bounded coefficients and $w$ has the same boundary sign,
then suitable multiples of $w$ plus $\varphi(x_0)\pm\varepsilon$ give a barrier.
The same construction with local boundary and outer-boundary bounds defines a
local barrier and yields the conclusion of Lemma 6.12 for bounded solutions.

If $\Omega$ satisfies an exterior sphere condition at $x_0$, then
\begin{equation}\tag{6.45}
w(x)=\tau(R^{-\sigma}-|x-y|^{-\sigma})
\end{equation}
with suitable $\tau,\sigma>0$ satisfies $Lw\le-1$ near $x_0$ when $c\le0$ and
the coefficients are bounded.  Thus every boundary point of an exterior-sphere
domain admits a barrier.

\begin{theorem}[Continuous Dirichlet data under an exterior sphere condition]
\label{gt-ch06-thm-6-13}
\dcref{ch6:6.13}
\group{gt-ch06}
If $L$ is strictly elliptic in a bounded domain, $c\le0$, the coefficients and
$f$ are bounded and locally $C^\alpha$, and $\Omega$ satisfies an exterior
sphere condition, then every $\varphi\in C^0(\partial\Omega)$ admits a unique
solution $u\in C^0(\bar\Omega)\cap C^{2,\alpha}(\Omega)$ of the Dirichlet problem.
\end{theorem}

\begin{theorem}[Classical global Dirichlet solvability]
\label{gt-ch06-thm-6-14}
\dcref{ch6:6.14}
\group{gt-ch06}
Under Theorem 6.13, if $\Omega$ is $C^{2,\alpha}$ and
$f$, the coefficients, and $\varphi$ are respectively in
$C^\alpha(\bar\Omega)$, $C^\alpha(\bar\Omega)$, and $C^{2,\alpha}(\bar\Omega)$,
then the unique solution belongs to $C^{2,\alpha}(\bar\Omega)$.
\end{theorem}
\begin{proof}
\sketch Apply Theorem 6.13, straighten the boundary near each point, and use
Lemma 6.10 and the local boundary estimate to identify the solution with the
local $C^{2,\alpha}$ Dirichlet solution.
\end{proof}

\begin{theorem}[Fredholm alternative for the Dirichlet problem]
\label{gt-ch06-thm-6-15}
\dcref{ch6:6.15}
\group{gt-ch06}
Let $L$ be strictly elliptic with $C^\alpha(\bar\Omega)$ coefficients on a
$C^{2,\alpha}$ domain.  Either the homogeneous Dirichlet problem has only the
zero solution, in which case the inhomogeneous problem has a unique
$C^{2,\alpha}(\bar\Omega)$ solution for every admissible $(f,\varphi)$, or its
homogeneous solution space is finite dimensional.
\end{theorem}
\begin{proof}
\sketch Restrict $L$ to $\mathbb B=\{u\in C^{2,\alpha}(\bar\Omega):u=0\text{ on }
\partial\Omega\}$.  For $\sigma\ge\sup c$, $L-\sigma$ is invertible by Theorem
6.14 and its inverse is compact on $C^\alpha(\bar\Omega)$.  The equation
\begin{equation}\tag{6.46}
u+\sigma(L-\sigma)^{-1}u=(L-\sigma)^{-1}f
\end{equation}
is therefore covered by the Fredholm alternative. Applying $L-\sigma$ gives
the equivalent boundary problem
\begin{equation}
\tag{6.47}
Lu=f,\qquad u\in\mathbb B.
\end{equation}
This yields the two cases.
\end{proof}

\section*{6.4. Interior and Boundary Regularity}

\begin{lemma}[Interior Schauder regularity]
\label{gt-ch06-lem-6-16}
\dcref{ch6:6.16}
\group{gt-ch06}
If $u\in C^2(\Omega)$ solves $Lu=f$ and $f$, $a^{ij}$, $b^i$, $c$ are
$C^\alpha(\Omega)$, then $u\in C^{2,\alpha}(\Omega)$.
\end{lemma}
\begin{proof}
\sketch On each $B\Subset\Omega$, solve the constant-principal-part problem
\begin{equation}\tag{6.48}
a^{ij}D_{ij}v+b^iD_iv=f-cu\quad\text{in }B,\qquad v=u\quad\text{on }\partial B.
\end{equation}
Lemma 6.10 gives $v\in C^{2,\alpha}(B)$, and uniqueness gives $v=u$.
\end{proof}

\begin{theorem}[Higher interior regularity]
\label{gt-ch06-thm-6-17}
\dcref{ch6:6.17}
\group{gt-ch06}
If $u\in C^2(\Omega)$ solves $Lu=f$ with coefficients and $f$ in
$C^{k,\alpha}(\Omega)$, then $u\in C^{k+2,\alpha}(\Omega)$.  Smooth data give
$u\in C^\infty(\Omega)$.
\end{theorem}
\begin{proof}
\sketch Take difference quotients.  For a coordinate direction and step $h$,
\begin{equation}\tag{6.49}
L(\Delta_l^h u)=\Delta_l^h f-(\Delta_l^h a^{ij})D_{ij}u(\cdot+he_l)
-(\Delta_l^h b^i)D_i u(\cdot+he_l)-(\Delta_l^h c)u(\cdot+he_l).
\end{equation}
Interior estimates give uniform bounds and compactness; passing to $h\to0$
proves the $k=1$ case.  Differentiate repeatedly and induct.
\end{proof}

\begin{lemma}[Regularity along a boundary portion]
\label{gt-ch06-lem-6-18}
\dcref{ch6:6.18}
\group{gt-ch06}
Let $T$ be a $C^{2,\alpha}$ boundary portion, $u\in C^0(\bar\Omega)\cap
C^2(\Omega)$ solve $Lu=f$, and assume $u=\varphi$ on $T$, with $\varphi$,
$f$, and the coefficients in $C^{2,\alpha}(\bar\Omega)$,
$C^\alpha(\bar\Omega)$, and $C^\alpha(\bar\Omega)$ respectively.  Then
$u\in C^{2,\alpha}(\Omega\cup T)$.
\end{lemma}
\begin{proof}
\sketch Cut off a small $C^{2,\alpha}$ subdomain whose boundary contains a
portion of $T$, extend its boundary data, approximate smoothly, and use the
Fredholm solvability, Corollary 6.7, and interior compactness.
\end{proof}

\begin{theorem}[Higher global regularity]
\label{gt-ch06-thm-6-19}
\dcref{ch6:6.19}
\group{gt-ch06}
If $\Omega$ is $C^{k+2,\alpha}$, $u\in C^0(\bar\Omega)\cap C^2(\Omega)$ solves
$Lu=f$, $u=\varphi$ on $\partial\Omega$, and the coefficients, $f$, and
$\varphi$ are in $C^{k,\alpha}(\bar\Omega)$, $C^{k,\alpha}(\bar\Omega)$, and
$C^{k+2,\alpha}(\bar\Omega)$, then $u\in C^{k+2,\alpha}(\bar\Omega)$.
\end{theorem}
\begin{proof}
\sketch Lemma 6.18 is the base case.  Flatten the boundary, take tangential
difference quotients, and recover the normal second derivative from
$D_{nn}u=(f-(L-a^{nn}D_{nn})u)/a^{nn}$.  Induction gives the general case.
\end{proof}

\section*{6.5. An Alternative Approach}

\begin{lemma}[Weighted Schauder estimate]
\label{gt-ch06-lem-6-20}
\dcref{ch6:6.20}
\group{gt-ch06}
If $u\in C^{2,\alpha}(\Omega)$ solves $Lu=f$, (6.12)--(6.13) hold, and
$|u|_{0;\Omega}^{(-\beta)}$, $|f|_{0,\alpha;\Omega}^{(2-\beta)}$ are finite, then
\begin{equation}\tag{6.50}
|u|_{2,\alpha;\Omega}^{(-\beta)}
\le C\bigl(|u|_{0;\Omega}^{(-\beta)}+|f|_{0,\alpha;\Omega}^{(2-\beta)}\bigr),
\end{equation}
where $C=C(n,\alpha,\lambda,\Lambda,\beta)$.
\end{lemma}
\begin{proof}
\sketch Apply Theorem 6.2 on $B_{d_x/2}(x)$ and multiply by the appropriate
powers of $d_x$. This first gives
\begin{equation}
\tag{6.51}
|u|_{2;\Omega}^{(-\beta)}
\leq C\bigl(|u|_{0;\Omega}^{(-\beta)}
+|f|_{0,\alpha;\Omega}^{(2-\beta)}\bigr).
\end{equation}
For the H\"older term split into $|x-y|\le d_x/4$ and its
complement, then take the supremum.
\end{proof}

\begin{lemma}[Weighted maximum estimate on a ball]
\label{gt-ch06-lem-6-21}
\dcref{ch6:6.21}
\group{gt-ch06}
If $B=B_R(x_0)$, $c\le0$, the coefficients of $L$ are bounded by $\Lambda$,
$u\in C^0(\bar B)\cap C^2(B)$ solves $Lu=f$ with $u=0$ on $\partial B$, and
$0<\beta<1$, then
\begin{equation}\tag{6.52}
\sup_B d_x^{-\beta}|u(x)|\le C\sup_B d_x^{2-\beta}|f(x)|,
\end{equation}
where $C=C(\beta,n,R,\lambda,\Lambda)$.
\end{lemma}
\begin{proof}
\sketch Compare $u$ with a linear combination of
$w_1=(R^2-|x-x_0|^2)^\beta$ and an exponential barrier.  The first controls
the annulus and the second controls the centre; the maximum principle gives
\begin{equation}
\tag{6.53}
|u(x)|\leq Nw(x)\qquad(x\in B),
\end{equation}
for $N=\sup d^{2-\beta}|f|$, hence (6.52).
\end{proof}

\begin{theorem}[Weighted Dirichlet solvability on a ball]
\label{gt-ch06-thm-6-22}
\dcref{ch6:6.22}
\group{gt-ch06}
Let $B$ be a ball, $0<\beta<1$, $c\le0$, and assume the hypotheses of Lemma
6.21 and $|f|_{0,\alpha;B}^{(2-\beta)}<\infty$.  Then the zero-boundary
Dirichlet problem has a unique solution
$u\in C^0(\bar B)\cap C^{2,\alpha}(B)$ with
$|u|_{0;B}^{(-\beta)}<\infty$, satisfying (6.50).
\end{theorem}
\begin{proof}
\sketch Apply the method of continuity to $L_t=tL+(1-t)\Delta$ on the weighted
spaces $|u|_{2,\alpha}^{(-\beta)}<\infty$ and
$|f|_{0,\alpha}^{(2-\beta)}<\infty$.  Lemma 6.21 and (6.50) give a uniform
operator bound; the Poisson endpoint is Theorem 4.9.
\end{proof}

\begin{corollary}[Continuous boundary data on a ball]
\label{gt-ch06-cor-6-23}
\dcref{ch6:6.23}
\group{gt-ch06}
Under Theorem 6.22, every $\varphi\in C^0(\bar B)$ gives a unique solution
$u\in C^0(\bar B)\cap C^{2,\alpha}(B)$ of $Lu=f$ in $B$, $u=\varphi$ on
$\partial B$.
\end{corollary}
\begin{proof}
\sketch Approximate $\varphi$ uniformly by smooth data, solve after subtracting
the approximants, and use the maximum principle and Corollary 6.3 to pass to
the limit.
\end{proof}

\section*{6.6. Non-Uniformly Elliptic Equations}

The equation
\begin{equation}\tag{6.54}
u_{xx}+y^2u_{yy}=0
\end{equation}
on $0<x<\pi$, $0<y<Y$ illustrates failure of arbitrary Dirichlet data.  A
Fourier mode satisfies $y^2f_n''-n^2f_n=0$, with exponents
$\beta_n=(1+\sqrt{1+4n^2})/2>0$ and
$\gamma_n=(1-\sqrt{1+4n^2})/2<0$.  Boundedness at $y=0$ forces the negative
mode to vanish and therefore forces zero data on $y=0$.

If $f/\lambda$, $|b|/\lambda$, and the domain are bounded, the Perron process
still gives a bounded interior solution when the subfunction set is nonempty
and bounded above.  At an exterior sphere point $x_0$ assume additionally
$|A(x)\cdot(x-y)|\ge\delta>0$ nearby; then
$a^{ij}(x_i-y_i)(x_j-y_j)\ge\lambda'|x-y|^2$ there, and the barrier (6.45)
remains valid.  Alternatively, after normalizing $\lambda=1$, a strict
exterior plane condition gives the local barrier
\[
w(x)=e^{\gamma d}(1-e^{-\gamma x_1}),\qquad \gamma\ge1+\sup|b|/\lambda.
\]

\begin{theorem}[Dirichlet solvability with non-uniform ellipticity]
\label{gt-ch06-thm-6-24}
\dcref{ch6:6.24}
\group{gt-ch06}
Let $L$ be elliptic in a bounded domain, $c\le0$, with
$a^{ij},b^i,c,f\in C^\alpha(\Omega)$ and $b^i,c,f$ bounded.  If $\Omega$ has
an exterior sphere at every boundary point and a strict exterior plane at every
point where some $a^{ij}$ is unbounded, then every continuous boundary datum
has a unique solution in $C^0(\bar\Omega)\cap C^{2,\alpha}(\Omega)$.
\end{theorem}

\begin{corollary}[Strictly convex domain]
\label{gt-ch06-cor-6-24-prime}
\dcref{ch6:6.24'}
\group{gt-ch06}
If $a^{ij}\in C^\alpha(\Omega)$ and $\Omega$ is bounded and strictly convex,
then $a^{ij}D_{ij}u=0$ has a solution in
$C^0(\bar\Omega)\cap C^{2,\alpha}(\Omega)$ for every continuous boundary datum.
\end{corollary}

Let $d(x)=\operatorname{dist}(x,\partial\Omega)$ near a $C^2$ boundary point
and let $\psi\ge0$ vanish only at $x_0$.  A sufficient local-barrier condition is
\begin{equation}\tag{6.55}
K(a^{ij}D_{ij}d+b^iD_i d)\le-k|L\psi|-|f-c\varphi(x_0)|
\end{equation}
for suitable $K,k>0$.  If the principal curvatures with respect to the inner
normal are $\kappa_i$ and $a_i$ are the corresponding principal coefficients,
\begin{equation}\tag{6.56}
(a^{ij}D_{ij}d)(x_0)=-\sum_{i=1}^{n-1}a_i\kappa_i.
\end{equation}
Thus the sufficient condition is
\begin{equation}
\tag{6.57}
\sum_{i=1}^{n-1}a_i\kappa_i>\sup_\Omega|b|,
\end{equation}
and with continuous $b$ the sharper condition is
\begin{equation}\tag{6.58}
\sum_{i=1}^{n-1}a_i\kappa_i-b_\nu>0.
\end{equation}
Let
\[
\Sigma_1=\{x_0:(6.58)\text{ holds}\},\qquad
\Sigma_2=\{x_0:a^{ij}v_i v_j\ne0\text{ for a normal }v\}.
\]

\begin{theorem}[Boundary criterion for degenerate ellipticity]
\label{gt-ch06-thm-6-25}
\dcref{ch6:6.25}
\group{gt-ch06}
If $L$ is elliptic in a $C^2$ domain, $c\le0$, and
$a^{ij},b^i/\lambda,c,f/\lambda\in C^\alpha(\Omega)\cap C^0(\bar\Omega)$,
with $\Sigma_1\cup\Sigma_2=\partial\Omega$, then every
$\varphi\in C^0(\partial\Omega)$ gives a unique solution
$u\in C^2(\Omega)\cap C^0(\bar\Omega)$.
\end{theorem}

\subsection*{6.7. Other Boundary Conditions; the Oblique Derivative Problem}

On $\mathbb R^n_+$ consider
\begin{equation}\tag{6.60}
\mathcal Nu=au+\sum_{i=1}^n b_iD_i u=\varphi\quad\text{on }x_n=0,
\end{equation}
with $b_n\ne0$; normalize first by
\begin{equation}
\tag{6.61}
b_n>0,\qquad |b|=(|b_t|^2+b_n^2)^{1/2}=1.
\end{equation}
For $a\le0$ and
$n\ge3$, $y^*=(y',-y_n)$, the Green function is
\begin{equation}\tag{6.62}
G(x,y)=\Gamma(x-y)-\Gamma(x-y^*)-2b_n\int_0^\infty e^{as}D_n\Gamma(x-y^*+bs)\,ds,
\end{equation}
and
\begin{equation}\tag{6.63}
\mathcal N_xG(x,y)=0\quad(x_n=0).
\end{equation}
Writing $G=\Gamma(x-y)-\Gamma(x-y^*)+\Theta(x,y)$,
\[
\Theta(x,y)=|x-y^*|^{2-n}g\!\left(\frac{x-y^*}{|x-y^*|},|x-y^*|\right),
\]
the function $g$ is smooth and
\begin{equation}\label{eq:6.64}\tag{6.64}
D_{x_i}\Theta=-D_{y_i}\Theta\ (i<n),\qquad D_{x_n}\Theta=D_{y_n}\Theta,qquad
|D^\beta\Theta|\le C|x-y^*|^{2-n-|\beta|}.
\end{equation}

\begin{theorem}[Flat-boundary oblique derivative estimate]
\label{gt-ch06-thm-6-26}
\dcref{ch6:6.26}
\group{gt-ch06}
Let $B_1=B_R(x_0)$, $B_2=B_{2R}(x_0)$ in the closed half-space and
$B_i^+=B_i\cap\mathbb R^n_+$, $T=B_2\cap\{x_n=0\}$.  If
$\Delta u=f$ in $B_2^+$, $\mathcal Nu=\varphi$ on $T$, $a\le0$, $b_n>0$,
$f\in C^\alpha$, and $\varphi\in C^{1,\alpha}$, then
\begin{equation}\tag{6.65}
|u|_{2,\alpha;B_1^+}\le C\bigl(|u|_{0;B_2^+}+R|\varphi|_{1,\alpha;T}+R^2|f|_{0,\alpha;B_2^+}\bigr),
\end{equation}
where $C=C(n,\alpha,b_n/|b|)$.
\end{theorem}
\begin{proof}
\sketch Represent the zero-boundary-data solution with $G$.  The reflected
Newtonian part $w_1$ obeys
\begin{equation}
\tag{6.66}
|D^2w_1|_{0,\alpha;B_1^+}
\leq C|f|_{0,\alpha;B_2^+},
\end{equation}
and (6.64) gives the same bound for the $\Theta$-potential $w_2$. Hence
\begin{equation}
\tag{6.67}
|D^2(w_1+w_2)|_{0,\alpha;B_1^+}
\leq C|f|'_{0,\alpha;B_2^+}.
\end{equation}
A cutoff gives the estimate for $\varphi=0$.
For general data use
\begin{equation}\tag{6.68}
\psi(x)=b_n^{-1}x_n\int_{\mathbb R^{n-1}}\varphi(x'-x_ny')\eta(y')\,dy',
\end{equation}
so that
\begin{equation}\tag{6.69}
\mathcal N\psi=\varphi\quad\text{on }x_n=0,
\end{equation}
and
\begin{equation}\tag{6.70}
b_n|\psi|_{2;B_2^+}\le CR|\varphi|_{1,\alpha;T}.
\end{equation}
Apply the zero-data estimate to $u-\psi$.
\end{proof}

\begin{lemma}[Weighted oblique estimate on a half-space domain]
\label{gt-ch06-lem-6-27}
\dcref{ch6:6.27}
\group{gt-ch06}
Under the hypotheses of Theorem 6.26 on a bounded half-space domain,
\[
|u|^*_{2,\alpha;\Omega\cup T}
\le C\bigl(|u|_{0;\Omega}+|\varphi|_{1,\alpha;T}+|f|_{0,\alpha;\Omega}\bigr).
\]
\end{lemma}
\begin{lemma}[Constant-coefficient weighted oblique estimate]
\label{gt-ch06-lem-6-28}
\dcref{ch6:6.28}
\group{gt-ch06}
The same estimate holds for the constant-coefficient operator $L_0$ of Lemma
6.1, with $C=C(n,\alpha,\lambda,\Lambda,b_n/|b|,\operatorname{diam}\Omega)$:
\begin{equation}\tag{6.71}
|u|^*_{2,\alpha;\Omega\cup T}
\le C\bigl(|u|_{0;\Omega}+|\varphi|_{1,\alpha;T}+|f|_{0,\alpha;\Omega}\bigr).
\end{equation}
\end{lemma}

\begin{lemma}[Variable-coefficient flat-boundary oblique estimate]
\label{gt-ch06-lem-6-29}
\dcref{ch6:6.29}
\group{gt-ch06}
Let $\Omega\subset\mathbb R^n_+$ have flat boundary portion $T$ and let
$u\in C^{2,\alpha}(\Omega\cup T)$ solve $Lu=f$ with
\begin{equation}\tag{6.72}
\mathcal N(x')u=\gamma(x')u+\sum_i\beta_i(x')D_i u=\varphi(x')\quad(x'\in T),
\end{equation}
where $|\beta_n|\ge\kappa>0$, $a^{ij},b^i,c,f\in C^\alpha$, and
$\gamma,\beta_i,\varphi\in C^{1,\alpha}$ with norms bounded by $\Lambda$.  Then
\begin{equation}\tag{6.73}
|u|^*_{2,\alpha;\Omega\cup T}
\le C\bigl(|u|_{0;\Omega}+|\varphi|_{1,\alpha;T}+|f|_{0,\alpha;\Omega}\bigr),
\end{equation}
where $C=C(n,\alpha,\lambda,\Lambda,\kappa,\operatorname{diam}\Omega)$.
\end{lemma}
\begin{proof}
\sketch Replace $u$ by $ue^{kx_n}$ to arrange $\gamma\le0$, $\beta_n>0$.
Freeze coefficients in a ball and rewrite the boundary condition as
\begin{equation}\tag{6.74}
\mathcal N(x_0')u=[\mathcal N(x_0)-\mathcal N(x')]u+\varphi\equiv\Phi.
\end{equation}
Lemma 6.28 gives
\begin{equation}\tag{6.75}
d_{x_0}^{2+\alpha}\frac{|D^2u(x_0)-D^2u(y_0)|}{|x_0-y_0|^\alpha}
\le\frac{C}{\mu^{2+\alpha}}\bigl(|u|_0+|\Phi|_{1,\alpha}+|F|_{0,\alpha}\bigr)
+\frac4{\mu^\alpha}[u]^*_{2;\Omega\cup T}.
\end{equation}
The coefficient-difference term in $\Phi$ is bounded exactly as in Theorem
6.2, and choosing $\mu$ small absorbs the seminorm term.
\end{proof}

\begin{theorem}[Global oblique derivative Schauder estimate]
\label{gt-ch06-thm-6-30}
\dcref{ch6:6.30}
\group{gt-ch06}
Let $\Omega$ be $C^{2,\alpha}$, $u\in C^{2,\alpha}(\bar\Omega)$ solve $Lu=f$, and
\[
\mathcal N(x)u=\gamma(x)u+\beta_i(x)D_i u=\varphi(x)\quad(x\in\partial\Omega),
\]
with
\begin{equation}
\tag{6.76}
|\beta_\nu|\geq\kappa>0\qquad\text{on }\partial\Omega,
\end{equation}
$f\in C^\alpha$, $\varphi\in C^{1,\alpha}$,
$a^{ij},b^i,c\in C^\alpha$, and $\gamma,\beta_i\in C^{1,\alpha}$.  Then
\begin{equation}\tag{6.77}
|u|_{2,\alpha;\Omega}\le C\bigl(|u|_{0;\Omega}+|\varphi|_{1,\alpha;\Omega}+|f|_{0,\alpha;\Omega}\bigr),
\end{equation}
where $C=C(n,\alpha,\lambda,\Lambda,\kappa,\Omega)$.
\end{theorem}
\begin{proof}
\sketch Cover the boundary by flattened charts, apply Lemma 6.29, and combine
with interior estimates as in Theorem 6.6.
\end{proof}

\begin{theorem}[Oblique derivative solvability]
\label{gt-ch06-thm-6-31}
\dcref{ch6:6.31}
\group{gt-ch06}
Assume $L$ is strictly elliptic with $c\le0$ and $C^\alpha(\bar\Omega)$
coefficients on a $C^{2,\alpha}$ domain, and let
$\mathcal N=\gamma+\beta\cdot D$ with $\gamma(\beta\cdot\nu)>0$ on $\partial\Omega$,
$\gamma,\beta\in C^{1,\alpha}$.  Then
\begin{equation}\tag{6.78}
Lu=f\ \text{in }\Omega,\qquad \mathcal Nu=\varphi\ \text{on }\partial\Omega
\end{equation}
has a unique $C^{2,\alpha}(\bar\Omega)$ solution for every
$f\in C^\alpha(\bar\Omega)$ and $\varphi\in C^{1,\alpha}(\partial\Omega)$.
\end{theorem}
\begin{proof}
\sketch Join $(L,\mathcal N)$ to $(\Delta,\partial_\nu+I)$ by
\begin{equation}\tag{6.79}
L_t=tL+(1-t)\Delta,\qquad \mathcal N_t=t\mathcal N+(1-t)(\partial_\nu+I).
\end{equation}
Uniform obliqueness and Theorem 6.30 give a $t$-independent Schauder bound.
After dividing by a positive barrier $\omega=c_1-c_2e^{\mu x_1}$, the maximum
principle bounds $|u|_0$ by $C(|f|_0+|\varphi|_0)$; hence
\begin{equation}\tag{6.80}
|u|_{2,\alpha}\le C\bigl(|f|_{0,\alpha}+|\varphi|_{1,\alpha}\bigr).
\end{equation}
Equivalently, for the operator
$\mathcal L_t=(L_t,\mathcal N_t):\mathscr B_1\to\mathscr B_2$,
\begin{equation}
\tag{6.81}
\|u\|_{\mathscr B_1}\leq C\|\mathcal L_tu\|_{\mathscr B_2},
\end{equation}
uniformly in $t$.
The method of continuity and solvability of the endpoint problem prove the claim.
\end{proof}

\section*{6.8. Appendix 1: Interpolation Inequalities}

\begin{lemma}[Interior interpolation inequality]
\label{gt-ch06-lem-6-32}
\dcref{ch6:6.32}
\group{gt-ch06}
If $j+\beta<k+\alpha$, $u\in C^{k,\alpha}(\Omega)$, and $\Omega$ is open, then
for every $\varepsilon>0$,
\begin{equation}\tag{6.82}
[u]^*_{j,\beta;\Omega}\le C(\varepsilon)|u|_{0;\Omega}+\varepsilon [u]^*_{k,\alpha;\Omega},
\end{equation}
and
\[
|u|^*_{j,\beta;\Omega}\le C(\varepsilon)|u|_{0;\Omega}+\varepsilon [u]^*_{k,\alpha;\Omega}.
\]
\end{lemma}
\begin{proof}
\sketch The base case $j=1$, $k=2$, $\alpha=\beta=0$ is
\begin{equation}
\tag{6.83}
[u]^*_1\leq C(\varepsilon)|u|_0+\varepsilon[u]^*_2.
\end{equation}
For $d=\mu d_x$, use finite differences on $B_d(x)$ when the two
points are close and the trivial $2|u|_0$ bound otherwise.  Apply the mean-value
theorem to derivatives and choose $\mu$ so the high-order seminorm coefficient
is at most $\varepsilon$.  The cases $\alpha=0$ or $\beta=0$ follow by induction
on the derivative order.
\end{proof}

\begin{lemma}[Weighted H\"older compactness]
\label{gt-ch06-lem-6-33}
\dcref{ch6:6.33}
\group{gt-ch06}
If $\Omega$ is bounded and a set $S\subset C_*^{k,\alpha}(\Omega)$ is bounded
and equicontinuous on $\bar\Omega$, then $S$ is precompact in
$C_*^{j,\beta}(\Omega)$ whenever $j+\beta<k+\alpha$.
\end{lemma}
\begin{proof}
\sketch Use Arzela--Ascoli for uniform convergence and (6.82) for convergence
in the lower norm.
\end{proof}

\begin{lemma}[Flat-boundary interpolation inequality]
\label{gt-ch06-lem-6-34}
\dcref{ch6:6.34}
\group{gt-ch06}
If $\Omega\subset\mathbb R^n_+$ has flat boundary portion $T$ and
$u\in C^{k,\alpha}(\Omega\cup T)$, then
\begin{equation}
\tag{6.89}
[u]^*_{j,\beta;\Omega\cup T}
\leq C|u|_{0;\Omega}+\varepsilon[u]^*_{k,\alpha;\Omega\cup T},
\end{equation}
and the corresponding norm inequality holds.
\end{lemma}
\begin{proof}
\sketch The first new case is
\begin{equation}
\tag{6.90}
[u]^*_2\leq C(\varepsilon)|u|_0+\varepsilon[u]^*_{2,\alpha},
\qquad\alpha>0.
\end{equation}
Repeat the proof of Lemma 6.32 in balls meeting $T$, using truncated
balls and the mean-value theorem along tangential segments.
\end{proof}

\begin{lemma}[Global interpolation inequality]
\label{gt-ch06-lem-6-35}
\dcref{ch6:6.35}
\group{gt-ch06}
If $\Omega$ is a $C^{k,\alpha}$ domain, $j+\beta<k+\alpha$, and
$u\in C^{k,\alpha}(\bar\Omega)$, then
\begin{equation}\tag{6.91}
|u|_{j,\beta;\Omega}\le C(\varepsilon,j,k,\Omega)|u|_{0;\Omega}
+\varepsilon|u|_{k,\alpha;\Omega}.
\end{equation}
\end{lemma}
\begin{proof}
\sketch Flatten the boundary, apply Lemma 6.34, and transfer norms by (6.30).
On a smaller chart this gives
\begin{equation}
\tag{6.92}
|u|_{j,\beta;B''}
\leq C(\varepsilon)|u|_{0;\Omega}
+\varepsilon|u|_{k,\alpha;\Omega}.
\end{equation}
A finite boundary covering yields, for each boundary-centred ball $B$ of a
uniform radius,
\begin{equation}
\tag{6.93}
|u|_{j,\beta;B\cap\Omega}
\leq C|u|_{0;\Omega}+\varepsilon|u|_{k,\alpha;\Omega}.
\end{equation}
Combine this with the interior inequality.
\end{proof}

\begin{lemma}[Global H\"older compactness]
\label{gt-ch06-lem-6-36}
\dcref{ch6:6.36}
\group{gt-ch06}
If $\Omega$ is a $C^{k,\alpha}$ domain and $S$ is bounded in
$C^{k,\alpha}(\bar\Omega)$, then $S$ is precompact in
$C^{j,\beta}(\bar\Omega)$ for $j+\beta<k+\alpha$.
\end{lemma}
\begin{proof}
\sketch Combine Arzela--Ascoli with the global interpolation inequality (6.91).
\end{proof}

\section*{6.9. Appendix 2: Extension Lemmas}

A locally finite partition of unity subordinate to an open cover is used to
patch local extensions.

\begin{lemma}[Extension from a H\"older domain]
\label{gt-ch06-lem-6-37}
\dcref{ch6:6.37}
\group{gt-ch06}
Let $\Omega$ be a $C^{k,\alpha}$ domain, $k\ge1$, and let $\Omega'$ contain
$\bar\Omega$.  For $u\in C^{k,\alpha}(\bar\Omega)$ there is
$w\in C_0^{k,\alpha}(\Omega')$ with $w=u$ on $\Omega$ and
\begin{equation}\tag{6.94}
|w|_{k,\alpha;\Omega'}\le C(k,\Omega,\Omega')|u|_{k,\alpha;\Omega}.
\end{equation}
\end{lemma}
\begin{proof}
\sketch In flattened coordinates extend across $y_n=0$ by
\[
\widetilde u(y',y_n)=\sum_{i=1}^{k+1}c_i\widetilde u(y',-y_n/i),\qquad
\sum_{i=1}^{k+1}c_i(-1/i)^m=1\ (0\le m\le k).
\]
This preserves derivatives through order $k$ and the H\"older seminorm.  Patch
the local extensions with a finite partition of unity and multiply by a cutoff.
\end{proof}

\begin{lemma}[Extension of boundary data]
\label{gt-ch06-lem-6-38}
\dcref{ch6:6.38}
\group{gt-ch06}
Let $\Omega$ be $C^{k,\alpha}$, $k\ge1$, and $\Omega'$ contain $\bar\Omega$.
Every $\varphi\in C^{k,\alpha}(\partial\Omega)$ has an extension
$\Phi\in C_0^{k,\alpha}(\Omega')$ satisfying $\Phi=\varphi$ on $\partial\Omega$.
The same holds locally for a $C^{k,\alpha}$ boundary portion and for
$\varphi\in C^0(\partial\Omega)\cap C^{k,\alpha}(T)$, with a $C^0$ extension
globally and $C^{k,\alpha}$ regularity near $T$.
\end{lemma}
\begin{proof}
\sketch In flattened coordinates set $\widetilde\Phi(y',y_n)=\widetilde\varphi(y')$,
then patch the local constructions by a partition of unity as in Lemma 6.37.
\end{proof}
